\section{Background}
\label{sec:background}

\subsection{LLM agents are entering HPC}
\label{sec:background-hpc}

Recent systems work demonstrates that the connection between LLM agents and HPC resources is now a functioning integration path rather than a hypothetical one. Ma et al.\ connect LangChain/LangGraph agents to HPC resources through Parsl and demonstrate agent-driven molecular dynamics workflows on ALCF's Polaris system \cite{Ma2025HPCAgent}. Rosendo et al.\ describe a modular AI-agent architecture --- an LLM interface, multi-agent decision-making, programmable facility APIs, and provenance-aware infrastructure --- deployed for cross-facility autonomous experimentation between ORNL's Manufacturing Demonstration Facility and its Oak Ridge Leadership Computing Facility \cite{Rosendo2025ORNL}. Mondesire et al.\ report a multi-agent LLM system that autonomously builds and repairs over 200 HPC software packages, reading compiler errors and build logs and iteratively revising build scripts \cite{Mondesire2025PEARC} --- a workflow that, by construction, feeds tool and compiler output directly into an agent's decision loop. Dawson et al.'s LARA-HPC frames the open problem explicitly: agentic systems can automate scientific workflows on supercomputers, but correctness, reproducibility, and \emph{safe interaction with computational resources} remain unresolved deployment obstacles \cite{Dawson2026LARA}. And outside the research literature, HPC centers are already responding to organic adoption: Aalto's scientific computing group documents that users have started running Claude Code and OpenAI Codex against their Triton cluster for coding assistance and Slurm job management, and explicitly flags prompt injection, malicious skills, external repositories, and code/data confidentiality as live concerns for site operators, not hypothetical ones \cite{AaltoSciComp2026}.

Taken together, this literature establishes that agentic HPC is real and growing, but its focus is capability and architecture --- provenance, reproducibility, orchestration --- rather than adversarial robustness. None of these papers evaluates what happens when the content an agent reads during normal operation is adversarial.

\subsection{LLM-agent security: known attack mechanisms}
\label{sec:background-agentsec}

The mechanism side of this problem is well studied outside HPC. Greshake et al.\ established indirect prompt injection as a practical attack class: adversarial instructions embedded in content an LLM-integrated application retrieves --- not content the user typed --- can hijack the application's behavior, exfiltrate data, and persist across sessions \cite{Greshake2023PromptInjection}. Ruan et al.'s ToolEmu shows that tool-using agents fail in materially harmful ways at a non-trivial rate even under careful evaluation, using an LM-emulated sandbox to make this measurable without needing to instantiate every real tool \cite{Ruan2024ToolEmu}. Debenedetti et al.'s AgentDojo operationalizes this further into a reusable benchmark: 97 realistic tasks and 629 security test cases spanning email, banking, travel, and workspace domains, explicitly designed to measure \emph{utility under attack} jointly with attack success rather than either alone \cite{Debenedetti2024AgentDojo} --- a methodological choice we adopt for the HPC setting (Section~\ref{sec:benchmark}). More recent work extends the attack surface: Ferrag et al.\ catalog over thirty attack techniques across input manipulation, model compromise, system/privacy attacks, and inter-agent protocol exploits (MCP, ACP, ANP, A2A) \cite{Ferrag2025ProtocolExploits}, and Wang et al.'s systematization of prompt-injection threats highlights that many existing defenses and benchmarks implicitly assume an agent should distrust its environment --- an assumption that breaks down for agents, like HPC agents, that are \emph{authorized to rely on} runtime observations such as file contents, logs, and tool output as a normal part of the job \cite{Wang2026PILandscape}. Shi et al.\ show that the vulnerability extends specifically to tool selection: a malicious tool document, crafted without white-box access to the target model, can be optimized to make an agent consistently choose an attacker-controlled tool over legitimate alternatives \cite{Shi2025ToolHijacker} --- a mechanism that maps directly onto HPC module systems, shared scripts, and MCP-based tool servers (Section~\ref{sec:tool-poisoning}).

This literature gives HPC a vocabulary of mechanisms --- indirect injection, tool-output poisoning, tool-selection manipulation, protocol-level exploits --- but its benchmarks and case studies are built around web and enterprise semantics: an email inbox, a bank account, a travel itinerary. None of it encodes what it means for an agent to hold scheduler authority, to have legitimate access to more than one project while a task should touch only one, or to interact with filesystem state that persists across jobs, users, and other agents.

\subsection{HPC security: established controls and their assumptions}
\label{sec:background-hpcsec}

HPC security has, in parallel, matured its own answers to a different question. NIST SP 800-223 defines a zone-based reference architecture --- access, management, compute, and storage zones --- that standardizes how HPC sites reason about and communicate their security posture, and analyzes HPC-specific threats against it \cite{NIST2024SP800223}. Prout et al.\ describe enhanced user separation deployed at the MIT Lincoln Laboratory Supercomputing Center: enforced isolation of processes, filesystem access, network traffic, and accelerators, designed so that each of over a thousand users effectively experiences a personal HPC system \cite{Prout2024UserSeparation}. Alam et al.\ describe a federated single-sign-on and zero-trust co-design spanning AI and HPC digital research infrastructure, unifying multi-factor authentication and time-limited role-based access across heterogeneous facilities \cite{Alam2024ZeroTrust}.

These are all strong, deployed answers to \emph{who is allowed to do what}. They are not designed to, and do not, answer a different question: given that an action was performed by a correctly authenticated, correctly authorized account, was the \emph{intent} behind that action the user's, or was it injected by something the agent read along the way? A POSIX permission check, a scheduler accounting record, or a zero-trust access decision has no way to distinguish ``the user asked for this'' from ``the agent was told to do this by a poisoned log file the user never saw.'' That distinction is exactly what the hijacked authorized agent threat model is about.
